\section{What the contribution changes, and what it imports}\label{sec:dependencies}

\subsection{The two distinct calibration chains}
For the twelve-integer equations the proof has the following form:
\begin{align*}
 \exists\boldsymbol x\in\Z^{12}\;\primary(\boldsymbol x)=0
 &\ \Longleftrightarrow\ 
 \text{the fixed source-34 native instance has a witness}\\
 &\ \Longleftrightarrow\ \operatorname{Halt}(\mathrm{K215},\text{blank})\\
 &\ \Longleftrightarrow\ \exists p\;\operatorname{Proof}_{\ZFC}(p,\bot).
\end{align*}
The first arrow is the new arithmetic reduction plus the identified
signed native kernel. The later arrows are the repository's ordinary
compiler and exact-table machine calibration. The companion equation has
the same chain. The complete inherited directory is retained at
\path{integer12/inherited/root_fusion/}; in particular,
\path{provenance/CALIBRATION_STATUS.md} distinguishes the ordinary
mathematical claim from the encoded proof-assistant target.
The inherited proof sources include the signed native arithmetic,
the finite-window tableau decoder, the source-34 coefficient adapter,
and the exact K215 ordinary calibration. Their original paths,
hashes, notices, and historical dependency pins are preserved.

The direct equation instead has the chain
\begin{align*}
 \exists\boldsymbol x\in\Z^{1190}\;
 Q_{\mathrm{direct}}(\boldsymbol x)=0
 &\ \Longleftrightarrow\ 
 \text{the fixed 522-natural-coordinate equation has a witness}\\
 &\ \Longleftrightarrow\ C_{222}(1)\\
 &\ \Longleftrightarrow\ \exists p\;\operatorname{Proof}_{\ZFC}(p,\bot).
\end{align*}
The first arrow is established by the finite sign certificate and
reversible substitutions in Section~\ref{sec:direct}. The natural
equation's arithmetic compiler and numerical-calculus calibration are
imported with their ordinary proofs and source pins. This route makes no
use of a Turing-machine transition table or a Jones universal polynomial.

The displayed implications are used as ordinary mathematical theorems
formalizable in ZFC on its own natural and integer domains. Composing
and negating them proves \eqref{eq:target}. No argument assumes that every
ZFC theorem is true, or that a finite run has searched all proofs.

\subsection{The formalization boundary is material}
The supplied derivations are ordinary mathematical proofs, supplemented
by finite exact certificates. No new complete Lean or Rocq development
was built. In particular, an externally stated theorem about a machine
and a proof predicate is distinct from a derivation of a quoted
biconditional in the project's internal \texttt{SetTheory.Form} calculus.
The latter requires an additional syntactic translation and formal
composition for all internally coded witnesses.

The archived sources themselves record these distinctions. We preserve
them instead of treating compilation, finite simulations, or hashes as
a substitute for an unbounded theorem. If a reviewer declines to accept
an imported ordinary calibration until independently re-establishing it,
the newly checked conclusion is the exact equivalence with that
calibrated parent. The final $\Con(\ZFC)$ interpretation then retains the
explicit calibration dependency. This is the usual dependency structure
of a mathematical reduction, and it identifies the remaining work for a
fully kernel-checked deliverable.

Existing formal work on Diophantine complexity bounds illustrates a
different scale of machine-checked result~\cite{bayerdavid}. It supplies
context, not a theorem automatically validating the present literal
polynomials or their coefficient instantiations.

\section{A small additional improvement: twelve positive variables}
\label{sec:supplement}
The supplemental directory contains a separate optimization of the
repository's ordinary twelve-positive-variable polynomial. Its variables
remain positive, its exact degree remains $70$, and its zero set is
unchanged. Only the evaluation schedule changes from $100$ to $98$
binary operations.

In the parent's unchanged registers set
\[
 d=c^4,\qquad e=1-Q,\qquad v=V+n,\qquad
 E_{\mathrm a}=QV^2-1,\qquad h=ev^2+E_{\mathrm a},\qquad b=eQ.
\]
Then the following is an identity over the polynomial ring:
\begin{equation}\label{eq:polarized}
 (dh+b)^2-dQ(2ev)^2=(dh-b)^2+4dbE_{\mathrm a}.
\end{equation}
Indeed the difference between the two squared terms is $4dhb$ and
$h-E_{\mathrm a}=ev^2$. The right side reuses $E_{\mathrm a}$ and saves
two operations in the complete circuit.

This is a modest evaluator improvement. Norm polarization is already
used elsewhere in the repository. It is not a new representation
theorem, variable reduction, or degree result. Its complete authentic
coefficient tuple is the \emph{compact44} tuple with
$\beta=2^{590069}-1$, preserved separately under
\path{supplemental/ordinary98/native/}.
It must not be mixed with the source-34 tuple in
Section~\ref{sec:equation}.

\section{Assessment and concrete next research questions}
\label{sec:assessment}

The most transferable observation is the prime-valuation separator.
It replaces size comparisons between radicals by congruence information
that the certificate already carries. Once the least squared valuations
are separated, the argument is insensitive to how large the witnesses
are. It also handles an initially negative guard radicand. The use of
the classical quintic is useful because its quadratic remainder has
$B\equiv1\pmod3$; monotonicity is no longer needed to prove $B\ne0$.

The rational auxiliary root is a second independent improvement. Root
extraction first proves it rational, the adjacent-Pell descent proves it
integral, and the final equation carries only the denominator-cleared
norm. This moves an integrality proof into the arithmetic argument and
reduces the polynomial's degree.

The direct result suggests treating sign obligations as part of the
optimization problem. A natural coordinate whose sign follows from
existing equations needs no separate guard. The finite certificate
makes this saving reviewable, and signed monic projection preserves it
without imposing an incorrect positivity condition on eliminated
coordinates.

These observations were not found in the inspected parent constructions
in the forms used here. That is a repository-relative statement, not a
claim of worldwide priority. The classical carrier, norm invariance,
Legendre's theorem, CRT, Pell identities, and monic substitution are
established mathematics. Further novelty assessment should examine
earlier relation-combining and valuation-separation literature, including
the broader Diophantine complexity tradition~\cite{mr}.

Several specific questions remain open in this packet:
\begin{enumerate}
 \item Can a different congruence arrangement separate the same roots
 with a lower-degree carrier or a lower-degree scale $S$? Every proposed
 replacement must retain intersectivity, exclusion of zero-scale roots,
 and full residue interpolation.
 \item Can the auxiliary Pell denominator be shared more efficiently
 with the main congruence certificate? A shorter circuit alone is not a
 proof that its total degree or zero set is preserved.
 \item Can multi-row sign deductions or coordinate sign changes improve
 the direct quartic beyond the $234$ guards needed by the present
 single-row rule family? The restricted optimality certificate leaves
 these stronger methods entirely open.
 \item Can the arithmetic reductions be formalized against the exact
 internal ZFC syntax, including coefficient binding and inherited
 unbounded calibration? This would strengthen the verification status,
 even if no numerical metric improved.
\end{enumerate}
No prospective bound from these questions is included in the result
table.

\section{Reproduction and package map}\label{sec:reproduction}
Python 3.10 or later and its standard library suffice for the arithmetic
checks. A TeX distribution with the packages named in
\path{report/report.tex} builds the article. From the package root:
\begin{verbatim}
python3 verify.py
python3 -O verify.py
cd report
pdflatex -interaction=nonstopmode -halt-on-error report.tex
pdflatex -interaction=nonstopmode -halt-on-error report.tex
\end{verbatim}
The root verifier regenerates the two twelve-integer circuits and checks
their exact output; checks the full generic norm identity by an
independent coefficient calculation; checks the coefficient payload
hashes, circuit structure and arithmetic front ends; rebuilds the direct
equation and checks its full sign and projection certificates; and
checks the identical-polynomial supplemental optimization. The normal
and optimized runs both use explicit exceptions, not assertions that
can disappear under \texttt{-O}.

The release manifest binds the deliverable files. Checks generate
temporary receipts by default and do not rewrite the retained source
certificates. The archived reports record the completed runs. A compact
map is:
\begin{center}
\small
\begin{tabular}{@{}p{.30\linewidth}p{.64\linewidth}@{}}
\toprule
Path & Contents\\
\midrule
\path{report/} & Article sources, bounded coefficient appendix, and PDF.\\
\path{integer12/} & Two new circuits, builders, independent arithmetic
checks, new proof, baseline circuits, and inherited native dependencies.\\
\path{direct/} & Literal 1190-variable quartic, residuals, 1224-variable
checkpoint, sign/projection certificates, independent checker, and
complete synchronized parent.\\
\path{supplemental/ordinary98/} & Identical twelve-positive polynomial,
98-operation evaluator, proof, and its separate coefficient tuple.\\
\path{reviews/} & Independent assistant reviews of the new mathematics
and literal direct artifacts.\\
\path{verification/} & Release receipts and PDF layout review.\\
\path{SOURCE_PINS.json} & Repository commit and received archive hashes.\\
\path{MANIFEST.json} & Final deliverable file hashes.\\
\bottomrule
\end{tabular}
\end{center}

The primary equations are supplied as exact straight-line integer
circuits because writing their fully expanded coefficient lists would
obscure both the definition and the proof. A straight-line circuit with
fixed integer leaves defines a unique ordinary polynomial. The direct
quartic is supplied in both residual and completely expanded form.

\appendix
\section{Complete bounded coefficient recipe}\label{app:coeff}

The following is the full source-34 finite arithmetic recipe~\cite{coeff},
reproduced from the incoming root-fusion report with its licenses retained. The defining payload and source bindings are preserved under \path{integer12/inherited/root_fusion/coefficients/}.
All ranges are bounded constants. The ordered \(943\) tile records
\((x_{i0},x_{i1},x_{i2},x_{i3})\), in west, north, east, south order,
are supplied in
\path{integer12/inherited/root_fusion/coefficients/recipe/wang943_tiles.json}.
Their ordering is part of the definition: selector \(0\) is Start
and selector \(1\) is Forbidden.
The supplied \texttt{k215.tm} contains all \(430\) working actions
from which the source builder reconstructs the whole ordered list.
Its SHA-256 is
\begin{center}\small\texttt{
fdbcef0825886035c3091c8537a69535bc96f9e265f0ae45863543d40ae70090}.
\end{center}
Thus this finite table is delivered data, not an unspecified machine
or a free coefficient parameter.

Put
\[
b=121,\quad L_0=8801,\quad d=bL_0=1064921,\quad
\rho=2^{121},\quad B_0=2^{1064921}.
\]
For \(0\le s<4\), \(0\le j<10\), define
\[
e_s(j)=o_s+2j,\qquad(o_0,o_1,o_2,o_3)=(1020,1000,980,960),
\]
and \(w_j=2^{10+2j}\) for \(0\le j<44\). Let
\[
\begin{split}
\mathcal E={}&\{0,\ldots,942\}
 \cup\{e_s(j):0\le s<4,\ 0\le j<10\}\\
 &\cup\{1063,1064,1095,1096,2195,3296\}.
\end{split}
\]
It has \(989\) elements and maximum \(3296\).
For each tile \(0\le i<943\), set
\[
c_i=1+\sum_{s=0}^3\sum_{j=0}^9
       \operatorname{bit}_j(x_{is}+512)w_{10s+j}
       +\sum_{j=40}^{43}w_j .
\]
Here \(\operatorname{bit}_j(m)=\lfloor m/2^j\rfloor\bmod2\)
is used only to define the finitely many fixed integer coefficients.
It is not an operation on a quantified Diophantine variable.
The other coefficients are
\[
c_{e_s(j)}=w_{10s+j},\quad
(c_{1095},c_{1096},c_{2195},c_{3296})=(w_{40},w_{41},w_{42},w_{43}),
\quad c_{1063}=c_{1064}=0.
\]
Define
\begin{align*}
\mu&=1022+\sum_{j=0}^{43}w_j,\\
Z_{\rm pad}&=(2^{119}-1)\sum_{e=8792}^{8798}\rho^e
                +(2^{27}-1)\rho^{8799},\\
D_C&=\rho^{40}+\rho^{41}+\rho^{1100}+\rho^{2200}
       +\sum_{e\in\mathcal E}c_e(\rho^{4393-e}+\rho^{5495-e}),\\
D_R&=\rho,\\
M_C&=B_0-1-\sum_{\substack{e\in\mathcal E\\e\ne1}}\rho^e
                   -2\rho^{1063},\\
M_{F,\rm native}
 &=4+\mu(\rho^{4393}+\rho^{5495})
     +\sum_{j=0}^9(\rho^{1000+2j}+\rho^{1021+2j})
     +\rho^{2195}+\rho^{3296}+Z_{\rm pad}.
\end{align*}
The six actual fixed ports are precisely
\[
\boxed{\begin{gathered}
\beta=B_0-1,\quad K_0=D_C+B_0D_R,\quad M_C,\quad
M_F=M_{F,\rm native}+B_0-1,\\
A_{\rm cone}=2^{1113},\qquad B_{\rm cone}=2^{666104}.
\end{gathered}}
\]
All displayed exponents are fixed integers. In particular,
\(M_F\) already includes its shift; shifting it again would change
the equation and invalidate the binding.

\begin{table}[ht]\centering
\begin{tabular}{lrr}\toprule
Fixed numeral & Bit length & Minimal unsigned bytes\\\midrule
\(\beta\)&1\,064\,921&133\,116\\
\(K_0\)&1\,065\,043&133\,131\\
\(M_C\)&1\,064\,921&133\,116\\
\(M_F\)&1\,064\,922&133\,116\\
\(A_{\rm cone}\)&1\,114&140\\
\(B_{\rm cone}\)&666\,105&83\,264\\\midrule
Total&&615\,883\\\bottomrule
\end{tabular}
\caption{Exact fixed payload sizes before compression.
The archive stores deterministic gzip encodings of the minimal
unsigned little-endian bytes and binds their raw hashes.}
\end{table}

The source proof verifies the whole coefficient interface, not just
these sizes. It establishes local tile decoding, carry-free geometry,
whole-cell rotation, mask populations, positive canonical slack, and
an explicit rank controller. In particular
\[
\operatorname{pc}(M_C)+\operatorname{pc}(M_{F,\rm native})=d,\quad
M_C\equiv2\pmod4,\quad M_{F,\rm native}\equiv4\pmod8.
\]
The controller uses
\(d=11^2\cdot13\cdot677\), \(\lambda=2341\), and a proved Boolean
subset coverage lemma; it does not assume that \(d\) is a power of
five. Its construction supplies odd canonical exponents, odd cubic
scale, and the canonical parity promises in \eqref{eq:canonical}.
The complete inherited proof and literal builder are included,
with their internal source pins and licenses preserved. See
\path{integer12/inherited/root_fusion/provenance/COEFFICIENT_PROOF.md}
and
\path{integer12/inherited/root_fusion/coefficients/recipe/source34_builder.py}.
The twelve-variable transformation changes none of these six numerals.



\begin{thebibliography}{12}
\bibitem{hls}
Andrea M. Hyde, Paul D. Lee, and Blair K. Spearman.
\emph{Polynomials $(x^3-n)(x^2+3)$ Solvable Modulo Any Integer}.
American Mathematical Monthly \textbf{121}(4) (2014), 355--358.
\href{https://doi.org/10.4169/amer.math.monthly.121.04.355}
{doi:10.4169/amer.math.monthly.121.04.355}.
The $n=19$ case supplies the classical intersective carrier;
an elementary proof is given here.

\bibitem{mr}
Yuri V. Matiyasevich and Julia Robinson.
\emph{Reduction of an arbitrary Diophantine equation to one in 13 unknowns}.
Acta Arithmetica \textbf{27} (1975), 521--553.
\href{https://doi.org/10.4064/aa-27-1-521-553}
{doi:10.4064/aa-27-1-521-553}.

\bibitem{bayerdavid}
Jonas Bayer and Marco David.
\emph{A Formal Proof of Complexity Bounds on Diophantine Equations}.
In ITP 2025, LIPIcs \textbf{352}, 3:1--3:18.
\href{https://doi.org/10.4230/LIPIcs.ITP.2025.3}
{doi:10.4230/LIPIcs.ITP.2025.3}.

\bibitem{repo}
Vladimir Reshetnikov and ProveIt contributors.
\emph{ProveIt}, repository snapshot
\texttt{85ae4070019efd66ea3e58e9a67d287528094e8b}.
\url{https://github.com/VladimirReshetnikov/ProveIt}.
The supplied source pins record all selected incoming archives.

\bibitem{interpolation}
ProveIt incoming archive.
\emph{ConZFC Twelve Integer Interpolation},
\path{docs/incoming/ConZFC_Twelve_Integer_Interpolation_2026-10-24.zip},
at the snapshot in~\cite{repo}.
The ordinary proof \path{proofs/INTERPOLATION12.md} and independent
audit are preserved under \path{integer12/proofs/}.

\bibitem{rootfusion}
ProveIt root-fusion package, inherited by~\cite{interpolation}.
Ordinary native/compiler proofs, normalized signed arithmetic,
coefficient bindings, and calibration status.
Preserved verbatim under \path{integer12/inherited/root_fusion/}.

\bibitem{coeff}
ProveIt source-34 coefficient construction.
\path{integer12/inherited/root_fusion/provenance/COEFFICIENT_PROOF.md},
with authentic native input files and ordinary source copies in the
same dependency tree. The finite recipe is reproduced in
Appendix~\ref{app:coeff}.

\bibitem{sync}
ProveIt incoming archive.
\emph{ConZFC Synchronized Certificates},
\path{docs/incoming/ConZFC_Synchronized_Certificates_2026-10-24.zip},
at the snapshot in~\cite{repo}. Complete received package preserved
under \path{direct/upstream/synchronized/}.

\bibitem{batched}
ProveIt incoming archive.
\emph{ConZFC Batched Certificates},
\path{docs/incoming/con_zfc_batched_certificates_2026-10-23.zip},
at the snapshot in~\cite{repo}. This supplies the 1690-integer comparison
point and the attributed signed monic projection implementation.

\bibitem{ordinary}
ProveIt ordinary-unit-products research directory, at the snapshot
in~\cite{repo}. Parent evaluator and ordinary proof preserved under
\path{supplemental/ordinary98/}, with the independently pinned compact44
native inputs and inherited notices.
\end{thebibliography}

